\chapter{Examen Blanc Chronométré (Conditions Réelles)}

\begin{tcolorbox}[colback=red!4!white,colframe=red!60!black,title=\textbf{Consignes de l'examen blanc}]
\begin{itemize}
    \item Reproduire la journée du 27 septembre : \textbf{Épreuve A --- Spécialité (3 heures)}, pause, puis \textbf{Épreuve B --- Rédaction (2 heures)}.
    \item Répondre sur feuille, \textbf{sans consulter les corrigés}, placés sur des pages séparées.
    \item Épreuve A : partie 1 = 30 QCM (une seule bonne réponse, 0,5 point chacun, 15 points) ; partie 2 = 3 exercices rédigés (25 points). Total sur 40, ramené sur 20.
    \item Épreuve B : sujet unique, noté sur 20 avec la grille du chapitre 1.
    \item Objectif de sécurité : au moins 12/20 à chaque épreuve.
\end{itemize}
\end{tcolorbox}

% =========================================================================
\section{Épreuve A --- Spécialité Informatique et Réseaux (3 heures, coefficient 4)}
% =========================================================================

\subsection*{Partie 1 : 30 QCM (45 minutes)}

\begin{enumerate}[leftmargin=*, label=\textbf{Q\arabic*.}]
    \item Un réseau /27 offre combien d'adresses hôtes utilisables ?
    \newline (A) 14 \quad (B) 30 \quad (C) 32 \quad (D) 62
    \item À quelle couche du modèle OSI fonctionne un commutateur (switch) classique ?
    \newline (A) 1 \quad (B) 2 \quad (C) 3 \quad (D) 4
    \item Quel protocole évite les boucles de commutation ?
    \newline (A) OSPF \quad (B) STP \quad (C) DHCP \quad (D) SNMP
    \item Quel protocole de routage est utilisé entre systèmes autonomes sur Internet ?
    \newline (A) RIP \quad (B) OSPF \quad (C) BGP \quad (D) EIGRP
    \item Dans IPsec, quel protocole assure le chiffrement des données ?
    \newline (A) AH \quad (B) ESP \quad (C) IKE \quad (D) GRE
    \item Quelle norme permet le transport de plusieurs VLAN sur un même lien ?
    \newline (A) 802.11ax \quad (B) 802.1X \quad (C) 802.1Q \quad (D) 802.3
    \item Quel mécanisme authentifie individuellement un poste avant l'accès au réseau filaire ou Wi-Fi ?
    \newline (A) NAT \quad (B) 802.1X avec RADIUS \quad (C) WEP \quad (D) DNS
    \item Quel est le port par défaut de SSH ?
    \newline (A) 21 \quad (B) 22 \quad (C) 23 \quad (D) 25
    \item Quelle adresse est une adresse IPv4 privée ?
    \newline (A) 8.8.8.8 \quad (B) 172.20.5.1 \quad (C) 41.140.0.1 \quad (D) 200.1.1.1
    \item Pour la voix sur IP, quel protocole établit et termine les appels ?
    \newline (A) RTP \quad (B) SIP \quad (C) SMTP \quad (D) FTP
    \item Quelle fibre optique convient aux liaisons de plusieurs dizaines de kilomètres ?
    \newline (A) Multimode \quad (B) Monomode \quad (C) Cuivre catégorie 6 \quad (D) Coaxial
    \item Un SLA de 99,9 \% tolère une indisponibilité annuelle d'environ :
    \newline (A) 5 minutes \quad (B) 53 minutes \quad (C) 8 h 45 \quad (D) 3 jours 15 h
    \item Quel protocole sert à superviser les équipements réseau ?
    \newline (A) SNMP \quad (B) SIP \quad (C) LDAP \quad (D) NTP
    \item Quelle technologie opérateur fournit des réseaux privés virtuels multi-sites avec qualité de service garantie ?
    \newline (A) ADSL \quad (B) MPLS \quad (C) Wi-Fi 6 \quad (D) Bluetooth
    \item Que désigne le RTO dans un plan de reprise d'activité ?
    \newline (A) La perte de données maximale \quad (B) La durée maximale d'interruption \quad (C) Le coût du plan \quad (D) La fréquence des tests
    \item En SQL, quelle clause filtre les résultats d'une agrégation ?
    \newline (A) WHERE \quad (B) HAVING \quad (C) ORDER BY \quad (D) LIMIT
    \item Quelle forme normale exige que chaque attribut non-clé dépende de toute la clé primaire ?
    \newline (A) 1FN \quad (B) 2FN \quad (C) 3FN \quad (D) BCNF
    \item Une association « plusieurs à plusieurs » entre deux entités se traduit en relationnel par :
    \newline (A) une clé étrangère dans l'une des tables \quad (B) une table de jointure \quad (C) une vue \quad (D) un index
    \item Quel principe de la POO consiste à masquer les données derrière des méthodes ?
    \newline (A) Héritage \quad (B) Polymorphisme \quad (C) Encapsulation \quad (D) Surcharge
    \item Quel code HTTP signifie « ressource créée » ?
    \newline (A) 200 \quad (B) 201 \quad (C) 404 \quad (D) 500
    \item Dans une API REST, quelle méthode HTTP est utilisée pour supprimer une ressource ?
    \newline (A) GET \quad (B) POST \quad (C) PUT \quad (D) DELETE
    \item Quelle est la complexité moyenne du tri fusion ?
    \newline (A) $O(n)$ \quad (B) $O(n \log n)$ \quad (C) $O(n^2)$ \quad (D) $O(\log n)$
    \item Quelle structure de données fonctionne en FIFO ?
    \newline (A) Pile \quad (B) File \quad (C) Arbre binaire \quad (D) Table de hachage
    \item Que fait la commande Linux \texttt{chmod 750 script.sh} ?
    \newline (A) rwxr-x--- \quad (B) rw-r--r-- \quad (C) rwxrwxrwx \quad (D) r-x------
    \item Quel algorithme est asymétrique ?
    \newline (A) AES \quad (B) RSA \quad (C) SHA-256 \quad (D) 3DES
    \item Quelle attaque consiste à insérer du code SQL via un formulaire ?
    \newline (A) XSS \quad (B) CSRF \quad (C) Injection SQL \quad (D) DDoS
    \item Quel texte marocain régit la cybersécurité ?
    \newline (A) Loi 09-08 \quad (B) Loi 05-20 \quad (C) Loi 43-20 \quad (D) Loi 55-19
    \item Quelle norme définit un système de management de la sécurité de l'information ?
    \newline (A) ISO 9001 \quad (B) ISO/IEC 27001 \quad (C) ITIL \quad (D) COBIT
    \item Que signifie ETL ?
    \newline (A) Extract, Transform, Load \quad (B) Encrypt, Transfer, Log \quad (C) Edit, Test, Launch \quad (D) Export, Table, List
    \item Un conteneur Docker, par rapport à une machine virtuelle :
    \newline (A) embarque son propre noyau \quad (B) partage le noyau de l'hôte \quad (C) est toujours plus lent \quad (D) nécessite un hyperviseur de type 1
\end{enumerate}

\subsection*{Partie 2 : Exercices Rédigés (2 h 15)}

\begin{tcolorbox}[colback=blue!5!white,colframe=blue!75!black,title=\textbf{Exercice 1 --- Adressage et segmentation (9 points)}]
Le consulat général du Maroc dans une grande ville dispose du bloc \texttt{192.168.40.0/24}. Il doit accueillir : 50 postes de guichet et back-office, 20 téléphones IP, 6 serveurs locaux, 10 caméras et équipements techniques, un Wi-Fi visiteurs de 25 clients et une liaison point à point vers le routeur de secours 4G.
\begin{enumerate}
    \item Proposez un découpage VLSM (réseau, masque, plage, diffusion) en justifiant l'ordre d'allocation. (5 pts)
    \item Indiquez pour chaque sous-réseau le VLAN associé et les règles de filtrage principales. (2 pts)
    \item Le consulat souhaite passer en double pile IPv6 : donnez un exemple d'adressage IPv6 pour deux VLAN à partir du préfixe \texttt{2001:db8:40::/48} et expliquez l'intérêt. (2 pts)
\end{enumerate}
\end{tcolorbox}

\begin{tcolorbox}[colback=blue!5!white,colframe=blue!75!black,title=\textbf{Exercice 2 --- Base de données du suivi des demandes de visa (8 points)}]
On souhaite concevoir la base d'un système de traitement des demandes de visa électronique. Un demandeur (passeport unique, nationalité) dépose des demandes ; chaque demande concerne un type de visa (tourisme, affaires, études), est traitée par un agent d'un consulat ou du service central, et passe par des statuts datés (déposée, en instruction, acceptée, refusée, délivrée).
\begin{enumerate}
    \item Proposez le MCD (entités, associations, cardinalités). (3 pts)
    \item Déduisez le MLD. (2 pts)
    \item Écrivez la requête SQL donnant, pour chaque type de visa, le nombre de demandes acceptées et le délai moyen (en jours) entre le dépôt et l'acceptation pour l'année 2026. (3 pts)
\end{enumerate}
\end{tcolorbox}

\begin{tcolorbox}[colback=blue!5!white,colframe=blue!75!black,title=\textbf{Exercice 3 --- Sécurisation d'une plateforme exposée sur Internet (8 points)}]
La plateforme de rendez-vous consulaire subit des ralentissements et l'on soupçonne une exploitation automatisée des créneaux par des revendeurs, ainsi que des tentatives d'attaque.
\begin{enumerate}
    \item Décrivez l'architecture de sécurité que vous mettez en place devant et autour de l'application (au moins six composants ou mesures, justifiés). (4 pts)
    \item Proposez des mesures fonctionnelles et techniques contre la réservation automatisée et la revente de créneaux. (2 pts)
    \item Quelles obligations légales marocaines s'appliquent aux données traitées et à la gestion d'un incident ? (2 pts)
\end{enumerate}
\end{tcolorbox}

\newpage
% =========================================================================
\section{Corrigé de l'Épreuve A}
% =========================================================================

\subsection*{Corrigé des 30 QCM}
{\small
\begin{longtable}{|c|>{\raggedright\arraybackslash}p{4.6cm}|p{9.2cm}|}
\hline
\rowcolor{gray!20} \textbf{Q} & \textbf{Bonne réponse} & \textbf{Explication courte} \\ \hline
\endfirsthead
\hline
\rowcolor{gray!20} \textbf{Q} & \textbf{Bonne réponse} & \textbf{Explication courte} \\ \hline
\endhead
1 & B --- 30 & $2^{5} - 2 = 30$ (5 bits d'hôte). \\ \hline
2 & B --- couche 2 & Le switch commute selon les adresses MAC (Liaison). \\ \hline
3 & B --- STP & \textit{Spanning Tree Protocol} bloque les liens redondants pour éviter les boucles. \\ \hline
4 & C --- BGP & Protocole de routage inter-domaines d'Internet ; OSPF et RIP sont internes. \\ \hline
5 & B --- ESP & ESP chiffre et authentifie ; AH n'authentifie que ; IKE négocie les clés. \\ \hline
6 & C --- 802.1Q & Étiquetage des trames VLAN sur un trunk. \\ \hline
7 & B --- 802.1X avec RADIUS & Contrôle d'accès au port avec authentification individuelle. \\ \hline
8 & B --- 22 & 21 = FTP, 23 = Telnet (non chiffré), 25 = SMTP. \\ \hline
9 & B --- 172.20.5.1 & Plage privée 172.16.0.0/12 ; 41.x et 200.x sont publiques. \\ \hline
10 & B --- SIP & SIP signale les appels ; RTP transporte la voix. \\ \hline
11 & B --- monomode & Longues distances ; la multimode se limite à quelques centaines de mètres. \\ \hline
12 & C --- 8 h 45 & 0,1 \% de 8 760 heures $\approx$ 8,76 h. \\ \hline
13 & A --- SNMP & Protocole de supervision ; NTP synchronise l'heure. \\ \hline
14 & B --- MPLS & Commutation par étiquettes avec QoS garantie entre sites. \\ \hline
15 & B --- la durée maximale d'interruption & RTO ; la perte de données maximale est le RPO. \\ \hline
16 & B --- HAVING & Filtre après GROUP BY ; WHERE filtre avant. \\ \hline
17 & B --- 2FN & Élimination des dépendances partielles ; la 3FN traite les dépendances transitives. \\ \hline
18 & B --- une table de jointure & Elle contient les deux clés étrangères. \\ \hline
19 & C --- Encapsulation & Masquer l'état interne derrière des méthodes. \\ \hline
20 & B --- 201 & 200 = succès simple, 404 = introuvable, 500 = erreur serveur. \\ \hline
21 & D --- DELETE & GET lit, POST crée, PUT modifie. \\ \hline
22 & B --- $O(n \log n)$ & Tri fusion et tri rapide (en moyenne). \\ \hline
23 & B --- File & \textit{First In, First Out} ; la pile est LIFO. \\ \hline
24 & A --- rwxr-x--- & 7 = rwx (propriétaire), 5 = r-x (groupe), 0 = aucun droit (autres). \\ \hline
25 & B --- RSA & Paire clé publique / clé privée ; AES et 3DES sont symétriques, SHA-256 est un hachage. \\ \hline
26 & C --- Injection SQL & Parade : requêtes paramétrées. \\ \hline
27 & B --- Loi 05-20 & 09-08 = données personnelles, 43-20 = services de confiance, 55-19 = simplification administrative. \\ \hline
28 & B --- ISO/IEC 27001 & ITIL et COBIT sont des référentiels de gestion et de gouvernance, pas de sécurité. \\ \hline
29 & A --- Extract, Transform, Load & Chaîne d'alimentation d'un entrepôt de données. \\ \hline
30 & B --- partage le noyau de l'hôte & D'où sa légèreté par rapport à une VM. \\ \hline
\end{longtable}
}

\subsection*{Corrigé de l'Exercice 1 --- Adressage}
Bloc 192.168.40.0/24 (256 adresses). Allocation du plus grand au plus petit :
\begin{itemize}
    \item \textbf{VLAN 10 Postes} (50 $\rightarrow$ 64, /26) : 192.168.40.0/26, masque 255.255.255.192, plage .1 -- .62, diffusion .63.
    \item \textbf{VLAN 30 Wi-Fi visiteurs} (25 $\rightarrow$ 32, /27) : 192.168.40.64/27, masque 255.255.255.224, plage .65 -- .94, diffusion .95.
    \item \textbf{VLAN 40 Téléphonie} (20 $\rightarrow$ 32, /27) : 192.168.40.96/27, plage .97 -- .126, diffusion .127.
    \item \textbf{VLAN 60 Caméras et technique} (10 $\rightarrow$ 16, /28) : 192.168.40.128/28, masque 255.255.255.240, plage .129 -- .142, diffusion .143.
    \item \textbf{VLAN 50 Serveurs} (6 $\rightarrow$ 8, /29) : 192.168.40.144/29, masque 255.255.255.248, plage .145 -- .150, diffusion .151.
    \item \textbf{Liaison point à point} (/30) : 192.168.40.152/30, plage .153 -- .154, diffusion .155.
    \item Réserve : 192.168.40.156 -- .255.
\end{itemize}
Règles : visiteurs $\rightarrow$ Internet uniquement ; postes $\rightarrow$ serveurs sur ports applicatifs et tunnel vers Rabat ; caméras isolées, accessibles seulement depuis le serveur d'enregistrement ; téléphonie isolée avec QoS ; administration via bastion. IPv6 : à partir de 2001:db8:40::/48, un /64 par VLAN, par exemple 2001:db8:40:10::/64 pour les postes et 2001:db8:40:40::/64 pour la téléphonie ; intérêt : espace d'adressage illimité, fin du NAT, auto-configuration, préparation à la généralisation d'IPv6 chez les opérateurs et les partenaires.

\subsection*{Corrigé de l'Exercice 2 --- Base de données visa}
\textbf{MCD :} DEMANDEUR (id, n° passeport UNIQUE, nom, prénom, nationalité, date\_naissance) ; DEMANDE (id, date\_dépôt, motif, durée\_séjour) ; TYPE\_VISA (id, libellé, durée\_max, tarif) ; AGENT (id, matricule, nom) ; UNITE (id, libellé, type : consulat / central) ; STATUT\_DEMANDE (id, statut, date\_statut, commentaire). Associations : DEMANDEUR (1,n) dépose (1,1) DEMANDE ; DEMANDE (1,1) concerne (0,n) TYPE\_VISA ; DEMANDE (0,1) traitée\_par (0,n) AGENT ; AGENT (1,1) affecté\_à (1,n) UNITE ; DEMANDE (1,n) suit (1,1) STATUT\_DEMANDE (historique daté).

\textbf{MLD :} \texttt{demandeur(id\_demandeur PK, passeport UNIQUE, nom, prenom, nationalite, date\_naissance)} ; \texttt{type\_visa(id\_type PK, libelle, duree\_max, tarif)} ; \texttt{unite(id\_unite PK, libelle, type)} ; \texttt{agent(id\_agent PK, matricule, nom, id\_unite FK)} ; \texttt{demande(id\_demande PK, id\_demandeur FK, id\_type FK, id\_agent FK NULL, date\_depot, motif, duree\_sejour)} ; \texttt{statut\_demande(id\_statut PK, id\_demande FK, statut, date\_statut, commentaire)}.

\textbf{Requête :}
\begin{verbatim}
SELECT t.libelle,
       COUNT(*) AS nb_acceptees,
       AVG(s.date_statut - d.date_depot) AS delai_moyen_jours
FROM demande d
JOIN type_visa t ON t.id_type = d.id_type
JOIN statut_demande s ON s.id_demande = d.id_demande
WHERE s.statut = 'acceptee'
  AND s.date_statut BETWEEN '2026-01-01' AND '2026-12-31'
GROUP BY t.libelle
ORDER BY nb_acceptees DESC;
\end{verbatim}
(Selon le SGBD, la différence de dates s'écrit \texttt{DATEDIFF(s.date\_statut, d.date\_depot)} en MySQL.) Points bonus : index sur \texttt{statut\_demande(statut, date\_statut)}, contrainte d'unicité sur un statut « acceptée » par demande, données personnelles minimisées et durée de conservation définie.

\subsection*{Corrigé de l'Exercice 3 --- Sécurisation de la plateforme}
\begin{enumerate}
    \item \textbf{Architecture :} protection anti-DDoS et CDN devant le site ; reverse proxy et \textbf{WAF} filtrant les attaques web (OWASP Top 10) ; pare-feu nouvelle génération et IPS ; \textbf{DMZ} pour les serveurs frontaux, base de données dans une zone interne non exposée ; TLS partout avec certificats valides ; authentification forte des agents (MFA) et bastion d'administration ; journalisation centralisée vers un \textbf{SIEM} surveillé ; sauvegardes 3-2-1 et PRA testé ; tests d'intrusion réguliers ; durcissement des serveurs et correctifs.
    \item \textbf{Contre la réservation automatisée :} vérification de l'identité de l'usager (code envoyé par SMS ou e-mail, pièce d'identité unique par rendez-vous) ; CAPTCHA et limitation de débit par IP et par compte ; un seul rendez-vous actif par pièce d'identité et par service ; créneaux libérés de façon aléatoire et non annoncés ; détection des comportements automatisés (vitesse, empreinte du navigateur) ; nom du bénéficiaire imprimé sur la convocation et contrôlé au guichet ; canal prioritaire pour les cas urgents afin de réduire la pression.
    \item \textbf{Obligations légales :} loi 09-08 sur les données personnelles (déclaration à la CNDP, finalité, minimisation, sécurité, droits des personnes) ; loi 05-20 sur la cybersécurité et DNSSI (mesures minimales, déclaration des incidents à la DGSSI/maCERT) ; loi 43-20 si signature ou cachet électroniques ; obligations de secret professionnel des agents ; hébergement des données sur le territoire national conformément aux orientations de souveraineté numérique.
\end{enumerate}

\newpage
% =========================================================================
\section{Épreuve B --- Rédaction sur les Attributions du Ministère (2 heures, coefficient 2)}
% =========================================================================

\begin{tcolorbox}[colback=blue!5!white,colframe=blue!75!black,title=\textbf{Sujet inédit (examen blanc)}]
\textit{« Le Ministère des Affaires Étrangères, de la Coopération Africaine et des Marocains Résidant à l'Étranger a engagé en 2024 une réorganisation profonde de ses structures et de son réseau. En quoi la modernisation de l'outil diplomatique conditionne-t-elle la capacité du Royaume à défendre ses intérêts et à servir ses citoyens à l'étranger ? »}
\end{tcolorbox}

\textit{Traitez le sujet avant de lire la suite.}

\newpage
\section{Corrigé de l'Épreuve B}

\subsection*{1. Analyse du sujet (10 minutes)}
\begin{itemize}
    \item \textbf{Mots-clés :} « outil diplomatique » (réseau des postes, ressources humaines, formation, organisation, systèmes d'information), « modernisation » (décret 2.24.957, IMFRED, renouvellement des consuls, numérique, programme budgétaire « modernisation de l'outil diplomatique »), « défendre ses intérêts » (Sahara, diplomatie économique, influence), « servir ses citoyens à l'étranger » (MRE, services consulaires).
    \item \textbf{Piège :} rester sur l'organigramme sans montrer le lien entre moyens et résultats ; oublier le volet citoyens.
    \item \textbf{Angle de l'ingénieur :} la DSI, les réseaux et la donnée sont au cœur de l'outil diplomatique moderne.
\end{itemize}

\subsection*{2. Problématique et plan}
\textbf{Problématique :} dans quelle mesure la modernisation des structures, des compétences et des systèmes du ministère est-elle la condition d'une diplomatie efficace au service des intérêts du Royaume et de ses citoyens à l'étranger ?
\begin{itemize}
    \item \textbf{I. Un outil diplomatique étendu, engagé dans une modernisation structurelle}
    \begin{itemize}
        \item \textbf{A.} Un réseau et des missions étendus : 104 ambassades, 53 consulats, attributions multiples, budget en quatre programmes.
        \item \textbf{B.} La réorganisation de 2024 : pôles, nouvelles directions (DSI, diplomatie économique), IMFRED, renouvellement des consuls, numérique.
    \end{itemize}
    \item \textbf{II. Une modernisation qui conditionne l'efficacité de l'action extérieure}
    \begin{itemize}
        \item \textbf{A.} Défendre les intérêts : réactivité sur le dossier du Sahara, diplomatie économique outillée, sécurité de l'information.
        \item \textbf{B.} Servir les citoyens : consulats de nouvelle génération, refonte des institutions des MRE ; conditions : compétences, données, évaluation.
    \end{itemize}
\end{itemize}

\begin{copiemodele}{La modernisation de l'outil diplomatique}
\partiecopie{Introduction}
Le 21 novembre 2024, le décret n° 2.24.957 fixait les nouvelles attributions et l'organisation du ministère des Affaires étrangères, de la Coopération africaine et des Marocains résidant à l'étranger. Loin d'une simple réforme administrative, ce texte traduit une conviction : la diplomatie d'un pays vaut ce que vaut son outil. L'outil diplomatique désigne l'ensemble des moyens humains, organisationnels, matériels et technologiques par lesquels un État conduit son action extérieure : administration centrale, réseau des ambassades et consulats, formation des agents, systèmes d'information.

Le Maroc dispose de l'un des réseaux les plus étendus d'Afrique, avec environ 104 ambassades et missions permanentes et 53 consulats généraux, au service de missions multiples : défense de l'intégrité territoriale, coopération africaine, diplomatie économique et service de plus de cinq millions de Marocains du monde. La modernisation de cet outil constitue d'ailleurs le premier des quatre programmes du budget du ministère.

Dès lors, \textbf{dans quelle mesure la modernisation des structures, des compétences et des systèmes du ministère conditionne-t-elle la capacité du Royaume à défendre ses intérêts et à servir ses citoyens à l'étranger ?}

Nous présenterons d'abord un outil diplomatique étendu, engagé dans une modernisation structurelle (I), avant de montrer que cette modernisation conditionne directement l'efficacité de l'action extérieure (II).

\partiecopie{I. Un outil diplomatique étendu, engagé dans une modernisation structurelle}
\textit{Le ministère porte des missions et un réseau considérables (A), qu'une réorganisation profonde vient adapter aux exigences du temps (B).}

\textbf{A. Un réseau et des missions à la mesure des ambitions du Royaume}

Le ministère conduit la politique étrangère sous l'autorité de Sa Majesté le Roi, anime la coopération africaine, porte la diplomatie économique et assure la protection des Marocains du monde. Pour cela, il s'appuie sur un réseau présent sur tous les continents, dont la densité s'est accrue en Afrique et en Amérique latine au cours de la dernière décennie, et sur des institutions spécialisées comme l'Agence Marocaine de Coopération Internationale.

Ce réseau traite des flux considérables : correspondance diplomatique, millions d'actes consulaires, près de 386 000 visas électroniques délivrés entre 2022 et 2024, opération Marhaba qui accueille plus de quatre millions de personnes chaque été. Son budget, présenté au Parlement en quatre programmes, a permis la création de 155 postes en 2026, signe d'un effort soutenu.

\textbf{B. Une réorganisation profonde engagée en 2024}

Le décret 2.24.957 a regroupé les directions générales en pôles homogènes et complémentaires et créé de nouvelles directions, parmi lesquelles une Direction des Systèmes d'Information chargée des plans informatiques, de l'infrastructure et de la cybersécurité, et une Direction de la Diplomatie Économique. L'Académie Marocaine des Études Diplomatiques est devenue l'Institut Marocain de Formation, de Recherche et d'Études Diplomatiques, appelé à professionnaliser davantage la formation.

Cette modernisation touche aussi les personnes : en 2026, 35 \% des postes consulaires ont été renouvelés, avec une forte féminisation. Elle touche enfin les outils : portail consulat.ma, prise de rendez-vous en ligne, plateforme eConsulat, apostille électronique, bureau d'ordre digital et parapheur électronique.

\textit{Ces réformes ne sont pas une fin en soi ; leur justification réside dans les résultats qu'elles permettent d'obtenir sur les deux fronts de l'action extérieure.}

\partiecopie{II. Une modernisation qui conditionne l'efficacité de l'action extérieure}
\textit{Un outil moderne renforce la défense des intérêts du Royaume (A) et la qualité du service rendu à ses citoyens (B), à condition de réunir compétences, données et évaluation.}

\textbf{A. Défendre les intérêts du Royaume}

Sur le dossier du Sahara, la dynamique qui a conduit à la résolution 2797 du 31 octobre 2025 et aux pourparlers de 2026 exige des postes diplomatiques réactifs, informés en temps réel et capables de porter un message cohérent, comme les six principes énoncés à l'Assemblée générale le 22 septembre 2026. Cela suppose des communications sûres entre Rabat et les 104 ambassades, une veille structurée et une protection rigoureuse de l'information face aux cybermenaces, conformément à la loi 05-20 et à la Stratégie nationale de cybersécurité 2030.

La diplomatie économique, à l'approche du Mondial 2030, requiert des outils de veille, des bases de données d'opportunités et des tableaux de bord permettant d'évaluer l'action de chaque poste. Un ministère organisé en pôles, doté d'une direction dédiée et d'un institut de formation à l'économie, est mieux armé pour transformer les atouts du pays en investissements.

\textbf{B. Servir les citoyens à l'étranger}

Pour les Marocains du monde, la modernisation se mesure au guichet : délais d'attente, simplicité des démarches, continuité de service. Les consulats de nouvelle génération, connectés à l'état civil national et aux administrations partenaires, incarnent la diplomatie de proximité voulue par le Souverain. La refonte des institutions des MRE annoncée le 6 novembre 2024, avec une Fondation Mohammedia opérationnelle et un mécanisme de mobilisation des compétences, ne produira ses effets que si le réseau consulaire est capable de la relayer.

Trois conditions doivent être réunies. Les compétences d'abord : recruter des ingénieurs en réseaux, en développement et en cybersécurité, former les agents et fidéliser les talents. Les données ensuite : fiables, protégées selon la loi 09-08, interopérables, hébergées sur le territoire national. L'évaluation enfin : des indicateurs de performance par programme et par poste, dans la logique de résultats des finances publiques. L'ingénieur d'État est, sur ces trois points, un acteur central de la modernisation.

\partiecopie{Conclusion}
Le ministère dispose d'un outil diplomatique étendu, qu'il modernise en profondeur depuis 2024 dans ses structures, ses compétences et ses systèmes. Cette modernisation n'est pas un luxe : elle conditionne la réactivité sur les dossiers vitaux, l'efficacité de la diplomatie économique et la qualité du service rendu aux Marocains du monde.

À l'heure où l'intelligence artificielle et la guerre informationnelle transforment les relations internationales, la prochaine étape sera sans doute celle d'une diplomatie augmentée par la donnée : le Maroc saura-t-il en faire un levier de souveraineté ?
\end{copiemodele}

\subsection*{3. Grille d'auto-correction (sur 20)}
{\small
\begin{longtable}{|p{6.2cm}|c|p{8cm}|}
\hline
\rowcolor{gray!20} \textbf{Critère} & \textbf{Points} & \textbf{Ce qui est attendu} \\ \hline
Introduction complète & 3 & Accroche datée, définition de l'outil diplomatique, contexte, problématique, annonce du plan. \\ \hline
Plan et logique & 4 & Deux parties équilibrées reliant moyens et résultats, chapeaux et transitions. \\ \hline
Connaissance du ministère et faits & 5 & Décret 2.24.957, pôles, DSI, IMFRED, réseau, budget en quatre programmes, 155 postes, chiffres MRE et e-Visa, résolution 2797. \\ \hline
Vision de l'ingénieur d'État & 3 & Réseaux sûrs, données, cybersécurité, indicateurs : le numérique au service des missions. \\ \hline
Esprit critique et propositions & 3 & Conditions de réussite (compétences, données, évaluation), limites honnêtes. \\ \hline
Expression et ton diplomatique & 2 & Style sobre, neutralité, terminologie officielle, orthographe. \\ \hline
\textbf{Total} & \textbf{20} & \\ \hline
\end{longtable}
}

\subsection*{4. Erreurs qui font perdre des points}
\begin{itemize}
    \item Réciter un organigramme sans montrer à quoi servent les structures.
    \item Oublier l'un des deux volets du sujet (intérêts du Royaume / citoyens à l'étranger).
    \item Employer une terminologie contraire aux positions officielles sur le Sahara ou prendre parti sur l'actualité électorale.
    \item Inventer des intitulés de directions ou des chiffres ; préférer « selon le décret de 2024 » ou « selon les données officielles ».
\end{itemize}
